\chapter{Research Methodology}
\label{ch:methodology}


\section{Research Design}

This is a retrospective, observational prediction-modelling study on
de-identified critical-care data. It follows a pre-registered protocol
(\texttt{00\_preregistration.md}, locked before any modelling) and reports
according to TRIPOD+AI \cite{collins2024tripodai} and PROBAST+AI
\cite{moons2025probastai} guidelines. The central design innovation is
treating the label definition, validation split, and evaluation metric as
primary variables rather than fixed background choices.

All code is implemented in R (v4.x) and Python (v3.11), with a single
configuration module (\texttt{config.R}) as the sole source of truth for all
paths, clinical identifiers, and thresholds.

\paragraph{Literature search protocol.} The design decisions set out in this
chapter are warranted by evidence identified through a documented search, and
two of the contributions claimed in Section~\ref{sec:contrib} are claims about
what that search did \emph{not} find. The protocol (databases, date ranges,
search strings, inclusion and exclusion criteria, verification standard, and
yield) is specified in Table~\ref{tab:search_strategy} and the
identification-to-inclusion counts in Figure~\ref{fig:prisma}. It is a
structured narrative search rather than a systematic review: searches followed
a written protocol and every included source was verified against a primary
publisher, PubMed, PMC or repository record, but no title-and-abstract
screening pass was performed over a full de-duplicated yield and no second
reviewer screened independently. Every claim of the form ``no identified study
does $X$'' in this thesis is therefore stated as ``to the best of our
knowledge'' and is bounded by that coverage, which is recorded as a limitation
in Section~\ref{sec:limitations}.


\section{Data Sources}
\label{sec:data}

\subsection{Development Data: MIMIC-IV v3.1}

The primary dataset is MIMIC-IV v3.1 \cite{johnson2023mimiciv,johnson2023mimiciv_scidata,pollard2026physionet}, a large
de-identified critical-care database from Beth Israel Deaconess Medical Center
(BIDMC), Boston. It was accessed under a PhysioNet credentialing agreement. The
study uses eight clinical tables (Table~\ref{tab:data_sources}).

\begin{table}[H]
    \centering
    \begin{tabular}{lrr}
        \toprule
        \textbf{Table}     & \textbf{Exists} & \textbf{Size (MB)} \\
        \midrule
        admissions         & yes             & 19.9               \\
        patients           & yes             & 2.8                \\
        icustays           & yes             & 3.3                \\
        chartevents        & yes             & 3,502.6            \\
        labevents          & yes             & 2,593.0            \\
        prescriptions      & yes             & 606.7              \\
        microbiologyevents & yes             & 117.2              \\
        inputevents        & yes             & 401.3              \\
        \bottomrule
    \end{tabular}
    \caption{MIMIC-IV v3.1 source tables used in this study.}
    \label{tab:data_sources}
\end{table}

\subsection{External Validation Data: eICU-CRD v2.0}

External validation uses eICU-CRD v2.0 \cite{pollard2018eicu}, a
multi-centre database from over 200 US ICUs. It uses different electronic health
record systems and covers a broader range of hospitals and patient populations
than the single-centre MIMIC-IV, making it a meaningful external test.

\subsubsection{Transporting the label, not just the model}
\label{sec:label_transport}

External validation of a sepsis model requires transporting two things: the
model and the \emph{label}. The model transports trivially, it is a frozen
set of coefficients applied to harmonised features. The label does not, because
Sepsis-3 onset is defined on documentation artefacts (an antibiotic order, a
microbiology culture) whose recording conventions differ between databases far
more than physiology does. Three properties of eICU-CRD govern how the label is
constructed here.

\textbf{Drugs are recorded across two tables.} eICU-CRD has no item
dictionary; drugs are free-text \texttt{drugname} strings, and they are split
between \texttt{medication} (physician orders) and \texttt{infusionDrug}
(continuous infusions). Vasopressors and much of the intravenous antibiotic
exposure appear only in the latter. Both tables are read and the antibiotic
exposure is their union; vasopressor exposure is likewise taken from
\texttt{infusionDrug} so that the SOFA cardiovascular component is not fixed at
its no-pressor value.

\textbf{Microbiology coverage is the binding constraint.} The Sepsis-3
suspicion anchor requires a culture, and \texttt{microLab} covers
\pcExtMicroCoverage\% of the analysis cohort (\pcExtNMicroStays{} of
\pcExtNCohortStays{} unit stays). In the remaining stays the conjunction cannot
fire under \emph{any} variant, whatever the patient's physiology. Coverage of
the culture limb is an upper bound on the conjunction rather than a measure of
it: both limbs are documented for only \pcExtNBothLimbStays{} stays
(\pcExtBothLimbCoverage\%), and the counts narrow further once the window and
the organ-dysfunction criterion are applied
(Section~\ref{sec:external_coverage} states each stage separately). Including
those person-hours in the denominator would not estimate a low sepsis
incidence; it would report the coverage of one source table as though it were
an epidemiological quantity. The primary external analysis is therefore
restricted to the microbiology-covered sub-cohort, in which an onset is
observable in principle, and the reported external event rate is compared
directly against MIMIC-IV's on that basis.

\textbf{The organ-dysfunction criterion is a rise, not a level.} Onset is
gated on a SOFA increase of at least two points over a baseline, evaluated
inside each variant's own window and against each variant's own baseline
(admission for Variants A and B, rolling 24-hour for Variant C), exactly as in
the development cohort. An absolute SOFA threshold would be a different
criterion, and because missing components are scored at assumed-normal values
it would reject stays for incomplete documentation rather than for intact organ
function, which in a database with sparser laboratory coverage than MIMIC-IV
is a systematic bias towards the null.

Each of the three label variants is applied to eICU-CRD with its own
antibiotic--culture windows, so the external label sets differ from one another
as they do internally. This is verified rather than assumed: a
label-invariant comparator such as NEWS2 can only score an identical AUROC
across variants if the underlying label sets are identical, so the pipeline
asserts at run time that the three external onset sets are distinct and halts
if they are not.

A secondary \textbf{antibiotic-only} arm relaxes the culture requirement across
the full cohort and is reported separately. It measures how much of the event
count the culture requirement removes. Its target is not Sepsis-3, so it is
never pooled with the primary arm and never substitutes for it.


\section{Ethical Considerations}
\label{sec:ethics}

This study uses two publicly available, de-identified critical-care databases.
Neither contains protected health information as defined under HIPAA, and no
re-identification was attempted.

\textbf{MIMIC-IV v3.1} is distributed by the MIT Laboratory for Computational
Physiology via PhysioNet \cite{johnson2023mimiciv}. Access requires completion
of the Collaborative Institutional Training Initiative (CITI) ``Data or
Specimens Only Research'' course and signing a PhysioNet Credentialed Data Use
Agreement (DUA). The DUA prohibits re-identification and restricts data sharing
to credentialed collaborators. The principal investigator (the author) completed
CITI training and executed the DUA before any data access; credentials are on
file with PhysioNet.

\textbf{eICU-CRD v2.0} \cite{pollard2018eicu} is distributed under the same
PhysioNet credentialing framework and DUA terms.

Both databases were collected under institutional review board (IRB) approvals
held by the original data custodians (Beth Israel Deaconess Medical Center for
MIMIC-IV; Philips Healthcare for eICU-CRD), with a waiver of informed consent
for the de-identified research dataset. No additional IRB approval was required
for secondary analysis of these de-identified records, consistent with the US
Common Rule (45 CFR 46.104(d)(4)) and PhysioNet's data access policies.


\section{Cohort Construction}

The base cohort is built by applying three inclusion criteria uniformly across
all label variants:

\begin{itemize}
    \item Age $\geq 18$ at ICU admission.
    \item First ICU stay per patient only (to prevent within-patient correlation).
    \item ICU length of stay $\geq 4$ hours (to ensure enough data for hourly
    features and at least one prediction opportunity).
\end{itemize}

The resulting cohort contains \textbf{\pcNStays{} stays} in MIMIC-IV.
Figure~\ref{fig:consort} shows the patient flow through each exclusion step.

\begin{figure}[H]
    \centering
\adjustbox{max width=\textwidth}{%
    \begin{tikzpicture}[
        node distance=0.6cm,
        box/.style={rectangle, draw=ATUGreen, thick, rounded corners=2pt,
        fill=ATULightGreen!20, text width=7cm, align=center,
        font=\small, inner sep=6pt, minimum height=0.8cm},
        excl/.style={rectangle, draw=ATUOrange, thick, rounded corners=2pt,
        fill=ATUSand!40, text width=5cm, align=center,
        font=\small, inner sep=5pt, minimum height=0.7cm},
        arr/.style={-{Latex[length=2mm]}, thick, ATUGreen},
        exarr/.style={-{Latex[length=2mm]}, thick, ATUOrange},
    ]

\node[box] (total) {MIMIC-IV v3.1 ICU stays\\$n = \pcConsortNTotal$};
\node[excl, right=1.2cm of total] (e1) {Non-first ICU stay\\$n = \pcConsortNExclNotFirst$ excluded};
\draw[exarr] (total.east) -- (e1.west);

\node[box, below=of total] (first) {First ICU stay per patient\\(all ages $\geq 18$ in MIMIC-IV)\\$n = \pcConsortNFirstStay$};
\draw[arr] (total) -- (first);

\node[excl, right=1.2cm of first] (e2) {LOS $< 4$ hours\\$n = \pcConsortNExclLos$ excluded};
\draw[exarr] (first.east) -- (e2.west);

\node[box, below=of first] (final) {\textbf{Analysis cohort}\\$n = \pcConsortNFinal$};
\draw[arr] (first) -- (final);

\node[box, below left=0.8cm and -0.5cm of final, text width=3.2cm]
(va) {Variant A\\Sepsis: \pcNOnsetsA\\(\pcPctOnsetsA\%)};
\node[box, below=0.8cm of final, text width=3.2cm]
(vb) {Variant B\\Sepsis: \pcNOnsetsB\\(\pcPctOnsetsB\%)};
\node[box, below right=0.8cm and -0.5cm of final, text width=3.2cm]
(vc) {Variant C\\Sepsis: \pcNOnsetsC\\(\pcPctOnsetsC\%)};

\draw[arr] (final.south) -- ++(0,-0.3) -| (va.north);
\draw[arr] (final.south) -- (vb.north);
\draw[arr] (final.south) -- ++(0,-0.3) -| (vc.north);

    \end{tikzpicture}}
    \caption{CONSORT-style patient flow diagram. All \pcConsortNFinal{} stays enter
    all three label variants; only the sepsis outcome differs. Every count is
    generated from \texttt{02\_cohort\_flow}, written by script
    \texttt{02\_cohort\_labels.Rmd}. Sepsis counts are \textbf{cohort-wide},
    at any hour of the stay; the subset falling inside the 72-hour modelling
    horizon, which is what the models are trained and evaluated on, is reported
    in the results chapter.}
    \label{fig:consort}
\end{figure}


\section{Outcome Definition: Three Sepsis-3 Label Variants}
\label{sec:label_variants}

Following Cohen et al.\ \cite{cohen2024subtle} and Johnson et al.\
\cite{johnson2018comparative}, the Sepsis-3 onset label is defined three
different ways. All three are reasonable readings of the same consensus
definition \cite{singer2016sepsis3}. They differ in the
antibiotic-culture time windows and the SOFA baseline
(Table~\ref{tab:label_variants}).

\begin{table}[H]
    \centering\small
    \adjustbox{max width=\textwidth}{%
    \begin{tabular}{llccccc}
        \toprule
        \textbf{ID} & \textbf{Name} & \textbf{ABX$\to$cx} & \textbf{cx$\to$ABX}
        & \textbf{SOFA$-$} & \textbf{SOFA$+$} & \textbf{Baseline} \\
        \midrule
        A & Narrow      & 24\,h & 24\,h & 24\,h & 12\,h & Admission     \\
        B & Seymour-Std & 72\,h & 24\,h & 48\,h & 24\,h & Admission     \\
        C & Liberal     & 72\,h & 72\,h & 48\,h & 24\,h & Rolling 24\,h \\
        \bottomrule
    \end{tabular}}
    \caption{Three pre-registered Sepsis-3 label variants. ABX = antibiotic;
    cx = culture; SOFA$-$/SOFA$+$ = hours before/after suspicion time in which
    a SOFA rise of $\geq 2$ must be observed. Baseline = SOFA reference value.}
    \label{tab:label_variants}
\end{table}

In all three variants, a patient has sepsis if there is suspected infection
(antibiotic order paired with a culture within the specified windows) and the
SOFA score increases by at least 2 points. Surveillance cultures (e.g.\ MRSA
screens) are excluded.

Variant B (Seymour-Standard) is the \textbf{primary analysis}. Variants A and C
are pre-specified sensitivity analyses for the variance decomposition.


\section{Beyond Sepsis-3: A Treatment-Independent Label}
\label{sec:alt_labels}

Variants A, B and C are three readings of one definition, and all three fire on
the same trigger: a culture drawn near an antibiotic order. That shared anchor
is a record of a clinician's decision to treat, not a measurement of the
patient. Two consequences follow, and they must be separated.

The first is \textbf{definitional}. Sepsis-3 \cite{singer2016sepsis3} makes
onset a function of when treatment was given, so any variation in the
antibiotic--culture windows necessarily moves the onsets. Part of the
label-variant spread reported in Chapter~\ref{ch:evaluation} is therefore
guaranteed by the construction of the definition and would appear even if the
underlying physiology were identical across patients. The second is
\textbf{empirical}: whether a label that owes nothing to treatment timing picks
out different patients, at different hours, with different predictability. Only
the second is a claim about the world, and it cannot be settled inside a family
of labels that all share the anchor.

Sepsis-3 is a conjunction: \emph{suspected infection} \textbf{and}
\emph{acute organ dysfunction}. Under Protocol Amendment 2 the conjunction is
split and each limb is labelled on its own
(Table~\ref{tab:alt_label_variants}).

\begin{table}[H]
    \centering\small
    \begin{tabular}{llccp{3.1cm}}
        \toprule
        \textbf{ID} & \textbf{Name} & \textbf{Suspicion anchor}
        & \textbf{Organ dysf.} & \textbf{Uses treatment timing?} \\
        \midrule
        B & Seymour-Std    & required     & required     & \raggedright membership and timing\tabularnewline
        E & Anchor-Only    & required     & not required & \raggedright yes\tabularnewline
        D & Deterioration  & not required & required     & \raggedright \textbf{no}\tabularnewline
        F & Sepsis-3 Early & required     & required     & \raggedright membership only\tabularnewline
        \bottomrule
    \end{tabular}
    \caption{The two limbs of the Sepsis-3 conjunction, labelled separately
    (Protocol Amendment 2), and the standard onset timing applied to the intact
    conjunction (Variant F, Protocol Amendment 4). Variant E is Variant B's
    suspicion anchor with the SOFA criterion removed; Variant D is acute organ
    dysfunction with the infection criterion removed; Variant F has Variant B's
    membership exactly and differs only in when the onset is placed.}
    \label{tab:alt_label_variants}
\end{table}

\textbf{Variant E (Anchor-Only)} takes onset to be the suspicion time itself,
using Variant B's windows, with no organ-dysfunction requirement. Its label is
purely a record of clinician behaviour. Predicting it well means predicting who
gets cultured and started on antibiotics, which is the quantity a
treatment-anchored sepsis label partly encodes. The share of Variant B's
above-chance discrimination that Variant E reproduces,
\begin{equation}
    \phi_{\text{anchor}}
    = \frac{\mathrm{AUROC}_E - 0.5}{\mathrm{AUROC}_B - 0.5},
    \label{eq:anchor_fraction}
\end{equation}
is reported as the \textbf{anchor-attributable fraction}.

\textbf{Variant D (Deterioration)} takes onset to be the first hour after hour~1
at which a \emph{physiology-only} SOFA score rises at least 2 points above the
patient's admission baseline (the mean of each component over hours 0--3).
"Physiology-only" means the vasopressor and catecholamine terms are withheld,
so the cardiovascular component is scored from mean arterial pressure alone.
No antibiotic order, culture draw or drug administration enters the definition
at any point: Variant D is the treatment-independent labelling that the
pre-registered family lacked. Both variants are fitted with the identical model
specification, feature set and temporal split, so the label is the only thing
that changes.

The PhysioNet/CinC 2019 operationalisation \cite{reyna2020utility} was
considered as the \emph{treatment-independent} alternative label and rejected
for that purpose. It sets onset to
$\min(t_{\text{susp}}, t_{\text{SOFA}})$ with the same
culture-plus-antibiotic trigger, so membership is still anchored to treatment
timing and it would not have broken the circularity. The same objection applies
to SEP-1 and to the CDC Adult Sepsis Event \cite{dutta2026performance}, both of
which are defined on antibiotic days.

Two limits on interpretation are stated here and carried through the results.
Variant~E is not a sepsis model and its AUROC is not a sepsis result. Variant~D
is a threshold function of recorded physiology, and five of the six SOFA
components are themselves model inputs, so its discrimination is partly
self-fulfilling and its \emph{level} is not comparable with Variant~B's. What
is comparable, and what carries the empirical content, is whether Variant~D
selects different patients at different hours: reported as onset agreement
(Cohen's $\kappa$) and median onset shift in
Section~\ref{sec:label_anchor_results}.

\subsection{Standard Onset Timing (Variant F, Protocol Amendment 4)}
\label{sec:variant_f}

That rejection settles the circularity question and no more. The $\min$ rule
does not break the anchor, membership still requires the
antibiotic-culture pair, but it does change something the pre-registered
family holds fixed and never examines: \emph{when} the onset is placed. Anchored
\emph{membership} and anchored \emph{timing} are separable properties, and only
the second is what the treatment-anchored evaluation window tests.

Variants A--C assign onset $= t_{\text{susp}}$ whenever the SOFA criterion is
satisfied anywhere in the evaluation window, so the SOFA criterion gates
membership and contributes no timing information. Seymour et al.\ \cite{seymour2016assessment} and the PhysioNet/CinC 2019 Challenge
\cite{reyna2020utility} instead assign
\begin{equation}
    t_{\text{onset}} = \min\bigl(t_{\text{susp}},\, t_{\text{SOFA}}\bigr),
    \label{eq:early_onset}
\end{equation}
with $t_{\text{SOFA}}$ the first hour in the evaluation window at which the
$\geq 2$-point rise over the admission baseline is already present. Equation
\eqref{eq:early_onset} is the branch that permits an onset to precede the first
antibiotic or culture; dropping it makes the pre-treatment window empty as a
matter of arithmetic (Proposition~\ref{prop:h3} in
Chapter~\ref{ch:evaluation}).

\textbf{Variant F (Sepsis-3 Early-Onset)} restores it and changes nothing else.
The suspicion windows, SOFA threshold, evaluation window and baseline are
Variant B's, so the two variants label exactly the same stays; only the
timestamp moves, and only ever earlier. Variant F therefore isolates the timing
branch as cleanly as Variants D and E isolate the two limbs.

Three properties are worth stating in advance. First, the SOFA score used for
$t_{\text{SOFA}}$ is the \emph{full} score, vasopressor components included, as
Seymour et al.\ specify; vasopressor administration is itself a treatment, so an
onset placed by the cardiovascular component is free of the
antibiotic-and-culture anchor that defines the pre-treatment window but is not
free of treatment in general. Variant D remains the fully treatment-independent
comparison.

Second, F and B label the same stays septic, membership is identical at the
level of the stay, so the informative agreement statistic for F is the onset
\emph{shift}, not the overlap. Agreement is nonetheless not perfect once
restricted to the modelled horizon, which leads to the third property.

Third, moving an onset earlier can move it \emph{into} the 72-hour prediction
window. A stay whose $t_{\text{susp}}$ falls after hour 72 contributes no
labelled event to any model here, but if its SOFA rise occurred inside the
window, Variant F places the onset there and the stay becomes a modellable
event. This is a one-directional effect: $F \supseteq B$ within the horizon, and
the extra stays are reported in Section~\ref{sec:pretreatment_label}. It means
the choice of onset rule determines not only \emph{when} events are recorded but
\emph{how many} are visible to a real-time model at all.

Variant F is post-hoc, is never pooled with A--C, and takes no part in H1 or the
variance decomposition. Its contrast against B joins exploratory family E3,
which enlarges that family and makes every Benjamini--Hochberg adjustment within
it stricter than before the amendment.

Variants D, E and F are post-hoc and are never pooled with A--C. They take no
part in H1 or in the variance decomposition table, which are computed from
exactly the same inputs as before the amendments; their contrasts form their own
exploratory family (E3), corrected by Benjamini--Hochberg.


\section{Prediction Framing}
\label{sec:framing}

Following Lauritsen et al.\ \cite{lauritsen2021framing} and Reyna et al.\
\cite{reyna2020utility}:

\begin{itemize}
    \item \textbf{Observation window:} Expanding from ICU admission, uses all
    data available up to the current hour.
    \item \textbf{Estimand:} the \emph{cause-specific discrete-time hazard} of
    sepsis onset in the current hour, conditional on the stay being at risk at
    the start of it, $\Pr(\text{onset at hour } t \mid \text{at risk at } t,
    \mathbf{x}_t)$, with discharge and death as competing causes.
    \item \textbf{Prediction frequency:} Every hour.
    \item \textbf{Start time:} 1 hour after ICU admission.
    \item \textbf{End time:} 72 hours after ICU admission (right-censored).
    \item \textbf{Alerting window:} 6 hours. An alert counts as timely if it is
    issued in the six hours up to onset; this governs the Utility Score and the
    per-patient detection metrics, not the quantity the models estimate.
\end{itemize}

\paragraph{What is estimated and what is scored, stated separately.} These are
two different objects in this design and conflating them would misdescribe every
discrimination number in Chapter~\ref{ch:evaluation}. The models estimate a
per-hour hazard, and the ranking metrics (AUROC, AUPRC and the calibration
curves) are computed on that hazard against an indicator that takes the value
one at the onset hour and zero at every other at-risk hour. They therefore
measure how well a model separates the hour in which onset is recorded from the
hours that precede it in the same and other stays; they are not the AUROC of a
``sepsis within the next six hours'' classifier, and they should not be read as
one. The six-hour window enters through the clinically oriented metrics instead:
the Utility Score pays a reward that decays linearly in lead time over that
window (Section~\ref{sec:metrics}), episode sensitivity asks whether an
alert fired inside it, and the lead-time distribution is reported against it
directly.

The choice follows from the competing-risks formulation rather than sitting
beside it. A cumulative six-hour risk is a functional of the hazard together
with the competing-cause hazards over the same window, so estimating the hazard
is the more primitive operation and leaves the horizon a reporting decision
rather than a modelling one. The cost is that a single number cannot be quoted
as ``six-hour discrimination'', and the corresponding limitation is recorded in
Section~\ref{sec:limitations}.

The alerting window and the utility scoring rule match the PhysioNet/CinC 2019
Challenge \cite{reyna2020utility}; the estimand does not, the Challenge having
scored a binary label defined over the window itself.


\section{Person-Hour Table and Competing Events}
\label{sec:person_hours}

Each patient stay is expanded into a table with one row per person-hour at
risk. Following Heyard et al.\ \cite{heyard2019dynamic}, each hour is
assigned one of four outcomes:

\begin{itemize}
    \item 0: Still at risk (censored at this hour)
    \item 1: Sepsis onset (the event of interest)
    \item 2: Discharged alive without sepsis (competing event)
    \item 3: Died without sepsis (competing event)
\end{itemize}

Treating discharge and death as simple censoring would overstate sepsis
incidence and bias model estimates. The multinomial competing-risks formulation
avoids this problem.


\section{Feature Construction}
\label{sec:features}

Features are extracted at each person-hour from chartevents, labevents,
inputevents, and outputevents. The feature set includes:

\begin{enumerate}
    \item \textbf{Vital signs}, 6 (last value carried forward per hour): heart
    rate, respiratory rate, SpO$_2$, MAP (invasive or noninvasive), temperature
    ($^\circ$C), GCS total.
    \item \textbf{Laboratory values}, 5 (last value carried forward): lactate,
    creatinine, total bilirubin, platelet count, WBC.
    \item \textbf{Vasopressor indicators}, 2: any vasopressor (yes/no); any
    norepinephrine or epinephrine (yes/no).
    \item \textbf{Six-hour trends}, 3: rolling slopes for heart rate,
    respiratory rate and MAP.
    \item \textbf{Missingness indicators}, 6: whether each of lactate,
    creatinine, bilirubin, platelets, white cells and a blood gas
    (P/F ratio) was ordered in the current hour, one flag per test.
\end{enumerate}

Temperature in Fahrenheit is converted to Celsius. FiO$_2$ values above 1.0
(percentage format) are divided by 100.

Two boundaries of this list are worth stating because they are easy to
misread. The P/F ratio enters as an \emph{order indicator} only: it is used in
the SOFA respiratory component that constructs the label, and its measurement
flag is a covariate, but the ratio itself is not a model input, because it is
recorded for a minority of person-hours and its forward-filled value carries
more information about ventilation practice than about the current hour. And
there is no lactate slope: rolling slopes are computed for the three variables
sampled densely enough to support one.

The declared set therefore holds \textbf{22 covariates}, and the fitted model
carries \textbf{29 terms per outcome} once the six-column natural-spline
baseline in elapsed hours and the intercept are counted. The vector is defined
once, in \texttt{config.R}, and every stage that fits or transports a model
reads it from there; \texttt{require\_features()} raises an error rather than
proceeding if any declared predictor is absent from a panel.


\section{Model Classes}
\label{sec:models}

\subsection[Primary Model: Dynamic Discrete-Time Competing-Risks Hazard Model]%
{Primary Model: Dynamic Discrete-Time Compet\-ing-Risks Hazard Model}
\label{sec:primary_model_spec}

The primary model is a multinomial pooled logistic regression over person-hours,
following D'Agostino et al.\ \cite{dagostino1990pooled} and Heyard et al.\
\cite{heyard2019dynamic}:

\begin{equation}
    \lambda_r(t \mid \mathbf{X}_{it}) =
    \frac{\exp\!\bigl(\alpha_r(t) + \boldsymbol{\beta}_r^\top \mathbf{X}_{it}\bigr)}
    {1 + \sum_{s} \exp\!\bigl(\alpha_s(t) + \boldsymbol{\beta}_s^\top \mathbf{X}_{it}\bigr)}
    \label{eq:primary_model}
\end{equation}

where $r \in \{1=\text{sepsis},\, 2=\text{discharge},\, 3=\text{death}\}$.
The baseline hazard $\alpha_r(t)$ uses a restricted natural cubic spline with
5 knots, placed at quantiles of training-set hour within each variant.

This model predicts the risk of sepsis every hour. Under short intervals
and low per-interval event probabilities (both conditions met for hourly ICU
data with onset probability below 1\%), it is mathematically equivalent to
time-dependent Cox regression, without requiring the proportional-hazards
assumption.

\paragraph{Identification: the specification is penalised, and must be.}
The unpenalised version of Equation~\eqref{eq:primary_model} is
\textbf{completely separated} on this design. Driven to a tight tolerance by
iteratively reweighted least squares, the whole coefficient vector diverges
together, $|\beta|$ of order $10^{15}$ under \emph{every} label variant,
which is the signature of separation along a direction in covariate space rather
than of one badly behaved predictor; the design matrix is full rank, so this is
not aliasing. A separated likelihood has no finite maximiser, so the fitted
coefficients are not identified and the point at which a quasi-Newton optimiser
stops is arbitrary.

This is a trap worth naming, because the standard diagnostic does not catch it:
\texttt{nnet::multinom} returns \texttt{convergence = 0} for every variant here.
That flag reports only that BFGS's relative-change stopping rule fired; it says
nothing about whether an optimum exists. A ridge penalty of
$\lambda = 10^{-4}$ is therefore applied, which makes the penalised
log-likelihood strictly concave and the maximiser unique. The value is
\textbf{not tuned}: across a 20-point ridge path the temporal-test AUROC moves
by less than 0.02 for every variant, and that insensitivity, not an
optimisation, is the justification for the constant. Every fit is checked
after estimation and flagged if $\max|\beta|$ exceeds 20.

Penalisation is not a refinement here; it is a precondition. Coefficient
comparisons across label variants (H2, Section~\ref{sec:h2}) are only meaningful
against a unique maximiser, and an unpenalised fit can land in a degenerate
region under one label and a benign region under another, manufacturing an
apparent label effect where none exists.

\paragraph{Uncertainty for this model comes from the bootstrap, with a
cluster-robust covariance alongside it.}
Person-hours are nested within stays, so a cluster-robust sandwich estimator
clustered on stay is the natural choice for coefficient standard errors. It is
not available off the shelf: \texttt{sandwich} supplies a \emph{bread} method
for \texttt{multinom} objects but no score-function method, so the whole
sandwich family fails on this class. The obstacle is narrow rather than
fundamental, and the missing piece is supplied here. For row $i$ and
non-reference category $k$ the score with respect to the coefficient on
covariate $j$ is $x_{ij}(y_{ik} - p_{ik})$, a residual matrix the fitted object
already stores; providing it, and fitting with the Hessian retained, yields a
clustered covariance for the primary model.

Two qualifications attach to it and neither is cosmetic. Because the
specification is ridge-penalised out of necessity
(Section~\ref{sec:primary_model_spec}), what the estimator returns is the
covariance of the \emph{penalised} estimator: the bread is the inverse penalised
Hessian, so the two factors are mutually consistent, but this is a
ridge-regularised sandwich and is described as one throughout. And it gives
standard errors on \emph{coefficients}, not on H2's count of unstable
covariates, which is a threshold rule applied across three separate fits and
admits no analytic variance of its own. H2 therefore remains a descriptive
count (Section~\ref{sec:h2}).

What the covariance does permit is a diagnostic on that count, added as
\textbf{Protocol Amendment 7} (post-hoc, \S21 of
\texttt{00\_preregistration.md}). A rule that flags a coefficient for doubling
in magnitude cannot distinguish a coefficient that genuinely moves from one
estimated too imprecisely for a doubling to mean anything, so for each flagged
covariate the difference between its two extreme estimates is referred to the
pooled standard error of those two. The ratio is reported alongside the
pre-registered count and never in place of it
(Section~\ref{sec:h2_uncertainty}). It is deliberately not a test: the three
fits share stays, so the differenced estimates are correlated and a pooled
standard error assuming independence is conservative by an unquantified amount,
and the penalisation qualification above applies to it as well. No $p$-value is
computed for it and it enters no corrected family.

Every interval on a performance quantity in this thesis continues to come from
the stay-level cluster bootstrap of Section~\ref{sec:inference}, which
resamples whole stays and so respects the same clustering. The analytic
covariance is reported for the coefficients and does not enter any hypothesis
verdict.

\subsection{Comparator: Gradient-Boosted Trees (GBT)}

An XGBoost model is fitted under the same conditions for each label variant
and split. Settings: maximum depth 6, learning rate 0.1, early stopping on a
15\% internal validation set, log-loss objective.

\subsection{Comparator: Static Cox Model}

A standard Cox proportional-hazards model using only the first observation per
stay, included as a deliberately weak baseline to show the cost of ignoring
time-varying data. Schoenfeld residual testing is applied to verify the expected
violation of proportional hazards.

\subsection{Rule-Based Comparators}

NEWS2, qSOFA, and SIRS are computed at each person-hour from the same features.
NEWS2 is the headline comparator following Evans et al.\
\cite{evans2021ssc}.


\section{Validation Architecture}
\label{sec:validation}

\begin{table}[H]
    \centering\small
    \begin{tabular}{p{2.8cm}p{6.5cm}p{3.5cm}}
        \toprule
        \textbf{Split} & \textbf{How} & \textbf{Why} \\
        \midrule
        Internal (primary) & Time-based: train 2008--2016, test 2017--2019
        \cite{guo2022temporal} & Honest internal estimate \\
        Internal (secondary) & Random 75/25 patient-level split (seed = 42)
        & Measure split optimism (H4) \\
        Treatment-anchored & Hours before first culture or antibiotic
        \cite{kamran2024evaluation} & Treatment leakage (H3) \\
        External & eICU-CRD v2.0 \cite{pollard2018eicu}, no retraining
        & Generalisability (H5) \\
        \bottomrule
    \end{tabular}
    \caption{Validation architecture. All splits are patient-level: no patient
    appears in both training and test sets.}
    \label{tab:validation}
\end{table}


\section{Evaluation Metrics}
\label{sec:metrics}

\subsection{Primary Metrics}

\textbf{PhysioNet Utility Score} \cite{reyna2020utility} is the primary
metric. It rewards early true alerts and penalises late and false ones. AUROC
is reported only for comparison, because it hides the clinical failure mode
\cite{wang2025srpm}.

Each stay contributes a raw utility. A septic stay first alerted at hour $t$
within the reward window $[t_{\mathrm{onset}} - 6,\ t_{\mathrm{onset}}]$ earns
$(t_{\mathrm{onset}} - t)/6$, so a full-horizon warning earns 1 and an alert
that merely coincides with onset earns nothing; a septic stay alerted only in
$(t_{\mathrm{onset}},\ t_{\mathrm{onset}} + 3]$ earns $-0.05$ per late alert
hour; a septic stay never alerted in either window earns $-1$; and every alert
hour outside a reward window, on any stay, costs $-0.05$.

The raw total $U$ is then referenced to the two strategies that bound it. The
\emph{inaction} strategy never alerts and therefore earns $U_{\text{inact}} =
-1$ per septic stay; the \emph{optimal} strategy earns the full reward on every
septic stay and raises no false alarm, giving $U_{\text{opt}} = +1$ per septic
stay. The reported score is
\begin{equation}
    \widetilde{U} = \frac{U - U_{\text{inact}}}
                         {U_{\text{opt}} - U_{\text{inact}}},
    \label{eq:utility_norm}
\end{equation}
so that $\widetilde{U} = 1$ is a perfect forecaster, $\widetilde{U} = 0$ is
\emph{exactly} the no-alert strategy, and $\widetilde{U} < 0$ is worse than
issuing no alerts at all. The inaction baseline has to appear in both numerator
and denominator: normalising by the optimum alone would score a model that
raises no alert at $-1$ rather than at $0$, contradicting the meaning of the
metric. Equation~\eqref{eq:utility_norm} is verified against both bounds at run
time, a threshold above every predicted probability is required to return
exactly zero, and no operating point may exceed one.

\textbf{Operating points.} A score of this kind is a property of a model
\emph{and} a threshold. At an hourly event rate below 0.5 per cent no model's
predicted probability approaches the conventional 0.5 cut-off, so that cut-off
measures the no-alert strategy for a well-calibrated model and an arbitrary
point for a poorly calibrated one. The Utility Score is therefore reported
twice: at the 0.5 cut-off, for comparability with the literature, and at its
maximum over a threshold grid. The grid is defined on the \emph{alert rate}, the fraction of person-hours alerted, rather than on the probability scale,
spanning $10^{-5}$ to $0.5$ over sixty logarithmically spaced points, so that
it adapts to each model's calibration and covers every operating point a ward
could contemplate. The maximum is taken over thresholds that raise at least one
alert, and the alert burden, alert rate and per-stay detection rate that attain
it are reported alongside, because a utility maximum bought at an unusable
alarm volume is not a deployable operating point.

\textbf{Calibration} is assessed via calibration intercept and slope, following
Heyard et al.\ \cite{heyard2020evaluation}, and, following Van Calster et al.\
\cite{vancalster2019calibration}, is additionally reported as a flexible
calibration curve (restricted cubic spline of the outcome on the logit of
predicted risk) summarised by the Integrated Calibration Index (ICI), E50, E90,
E\textsubscript{max}, and the scaled Brier score. Predicted probabilities are
floored at $10^{-6}$ rather than at the conventional $10^{-3}$, because at an
hourly event rate below 0.5 per cent the conventional guard overwrites the
majority of the distribution being measured.

\subsection{Secondary Metrics}

AUPRC (important at low event rates), alert burden (false alerts per true alert;
benchmark: 1.4 from Moor et al.\ \cite{moor2023review}), lead time
(benchmark: 3.7 hours from Moor et al.), and decision-curve net benefit
\cite{vickers2006dca}. Net benefit is evaluated on a threshold grid running
from one tenth to ten times the observed event prevalence rather than on the
conventional 0.01--0.20 grid, which at this prevalence lies entirely above the
range of thresholds any deployment would use.

\subsection{Deployment-Facing Metrics}
\label{sec:deployment_metrics}

AUROC, AUPRC and the Utility Score characterise a model but not the alert load
a ward would carry. The following are therefore reported alongside them, at a
threshold set to the largest cut-off achieving a fixed per-hour sensitivity of
80 per cent rather than at an arbitrary 0.5 cut-off:

\begin{itemize}
    \item \textbf{Positive predictive value} and its reciprocal, the number
    needed to evaluate (NNE), at that fixed sensitivity.
    \item \textbf{False alerts per patient-day}, which normalises alarm volume
    by exposure time rather than by the number of true alerts.
    \item \textbf{Per-patient (event-level) detection}: the fraction of septic
    stays alerted within the actionable window
    $[\text{onset}-6\,\text{h},\ \text{onset}]$, the fraction alerted at any
    point before onset, patient-level PPV, and the fraction of all stays
    alerted. Per-hour and per-patient views of the same operating point are
    reported side by side, because the choice of denominator moves the headline
    PPV by more than an order of magnitude.
    \item \textbf{False-alarm episodes per patient-day}, counting a contiguous
    run of alert hours once rather than once per hour.
    \item \textbf{Lead time} from first alert to onset: median, interquartile
    range, and the fraction of septic stays alerted at least 3 and 6 hours
    ahead.
\end{itemize}

The same quantities are computed on eICU-CRD with thresholds re-derived on the
external cohort.


\section{Sample Size}
\label{sec:sample_size}

Sample size is calculated using the Riley et al.\ \cite{riley2019sample}
method, targeting a shrinkage factor $\geq 0.9$ with \pcNCandidatePredictors{}
parameters and an anticipated $C$-statistic of 0.846 from Moor et al.\
\cite{moor2023review}. Each label variant is evaluated separately; an
underpowered variant is a finding, not a failure.

\paragraph{The parameter count, and the criterion's fit to this model.}
\pcNCandidatePredictors{} is the number of terms the model estimates for one
outcome: 22 covariates, the six-column spline baseline and an intercept
(Section~\ref{sec:features}). The mapping is imperfect and the imperfection runs
against the study, so it is stated rather than left implicit. Riley et al.'s
criteria are derived for a \emph{binary} outcome, whereas the primary model is a
four-category multinomial estimating those \pcNCandidatePredictors{} terms for
\emph{each} of three non-reference outcomes, three times that count in all.
Taking the
per-outcome count treats each cause-specific sub-model as the unit, which is the
closest defensible reading of a binary criterion in this setting but still
understates the dimensionality of the whole fit. The minimum sample size scales
approximately linearly in the parameter count, so the margins reported below
absorb even the all-outcomes reading with a factor of ten to spare, and the
adequacy verdict does not depend on which reading is taken.

\paragraph{The unit the criterion applies to.} Riley et al.'s criteria are
derived for a sample of \emph{independent} observations, and person-hours are
not independent, they are nested within ICU stays at roughly 45 correlated
rows per stay, and one stay per patient. The adequacy calculation is therefore
parameterised on the \textbf{stay}: the outcome proportion supplied to the
criterion is the fraction of training stays labelled septic under that variant,
and the resulting minimum is compared against the number of training stays. The
person-hour count is reported alongside as the size of the design matrix and is
never used as a sample size. This is the same unit that governs every interval
in this thesis, for the same reason: the cluster bootstrap of
Section~\ref{sec:inference} resamples stays, not rows. Reporting adequacy
against 2.3 million person-hours would overstate the available information by
roughly the mean number of at-risk hours per stay, and would do so in the
direction that makes the study look better powered than it is.


\section{Statistical Inference and Multiplicity}
\label{sec:multiplicity}
\label{sec:inference}

\subsection{Uncertainty: the Stay-Level Cluster Bootstrap}

Person-hours are nested within ICU stays: a single stay contributes on average
about 45 highly correlated hourly rows. Treating those rows as independent
would produce standard errors several times too small. Every interval in this
thesis therefore comes from a \textbf{stay-level cluster bootstrap}, in which
whole ICU stays are drawn with replacement and all person-hours belonging to a
drawn stay are carried with it. Intervals are 95\% percentile intervals from
$B = 2{,}000$ replicates for internal and subgroup analyses and $B = 1{,}000$
for eICU-CRD, which contains 7.9 million person-hours. The seed is fixed at 42.

Contrasts between quantities measured on the \emph{same} person-hours are
evaluated within replicate, so the interval accounts for their correlation.
This applies to the label-variant comparison (all three variants are built on
the same test stays) and to the model-class comparison (all models score the
same hours). Contrasts across \emph{different} evaluation sets (temporal
versus random split, and MIMIC-IV versus eICU-CRD) use independent bootstrap
streams and propagate both variances.

For subgroup analyses the bootstrap is stratified: stays are resampled within
each subgroup so the observed subgroup sizes are held fixed. The subgroups are
the object of study, not a random draw from a larger frame.

P-values are obtained by recentring the bootstrap distribution on the null
boundary and reading off the tail area, using the $(1 + \text{count})/(B + 1)$
convention so that no p-value is reported as exactly zero. The smallest
reportable p-value is therefore $1/(B+1)$.

\subsection{Confirmatory and Exploratory Families}

The six pre-registered hypotheses form a single \textbf{confirmatory} family
and are corrected by Holm--Bonferroni at a family-wise error rate of 0.05. The
family size is fixed at $m = 6$: the number of hypotheses that were
pre-registered, not the number that turned out to be testable. H2 and H3 each
reserve a slot without contributing a p-value (Section~\ref{sec:decomposition}
explains why neither admits one), which is deliberately conservative. A
hypothesis becoming untestable after the data were seen must not be allowed to
relax the correction applied to the others.

All subgroup analyses are \textbf{exploratory}. The pre-registration fixed
their \emph{scope} (that race, language and insurance would be examined),
but not the subgroup levels, the reference level, or the direction of any
comparison. They are corrected by Benjamini--Hochberg at $q = 0.05$ within
three families: E1, all subgroup AUROC contrasts against the reference level
across every equity variable and both models; E2, all subgroup
label-sensitivity contrasts; and E3, the alternative-label contrasts of
Protocol Amendment 2, in which the AUROC under each limb of the Sepsis-3
conjunction (Section~\ref{sec:alt_labels}) is compared with Variant B for both
models. The anchor-attributable fraction of
Equation~\ref{eq:anchor_fraction} is a ratio with no meaningful null at zero
and is reported as an interval only, outside the corrected family. The
correction denominator is the number of
contrasts \emph{computed}, not the smaller number selected for presentation in
Chapter~\ref{ch:evaluation}. Bonferroni-adjusted p-values are also reported so
that a reader preferring family-wise control over false-discovery control can
substitute them.

The two procedures differ by design. A confirmatory claim must survive
family-wise control, because one false confirmation invalidates the claim. An
exploratory scan exists to generate hypotheses and needs only its
false-discovery proportion bounded; applying family-wise control to it would
suppress every signal the scan is meant to surface.

Each subgroup is contrasted against the largest level of its own variable. The
reference is chosen by sample size alone, never by outcome.

\paragraph{Two event floors, and what each decides (Protocol Amendment 6).}
Two distinct thresholds govern the subgroup analysis and the distinction is
load-bearing. \texttt{MIN\_SUBGROUP\_EVENTS} $= \pcMinSubgroupEventsEstimate$
decides whether a subgroup AUROC is \emph{computed}: below it, together with a
50 person-hour minimum, an AUROC is not estimable and the subgroup is listed as
such rather than shown as a point estimate.
\texttt{MIN\_SUBGROUP\_EVENTS\_INTERPRET} $= \pcMinSubgroupEvents$ decides
whether a computed AUROC is \emph{interpreted}, and whether its contrast enters
family E1. Strata below the interpretation floor are still estimated and still
tabulated, flagged as such, because the figure showing their intervals reaching
the chance line is itself the argument for the floor; they are excluded from
every statement about spread and from the corrected family. The constant is set
equal to the external-validation floor deliberately: it answers the same
question on the same warrant, so the two cannot drift apart.

This is \textbf{Protocol Amendment 6}, adopted post-hoc on 6 August 2026 and
recorded as \S20 of the pre-registration. It must be read against the rule
stated two paragraphs above, namely that the correction denominator is the
number of contrasts computed rather than the number selected for presentation.
Restricting E1 from thirty computed contrasts to \pcNSubgroupContrasts{} is on
its face the manoeuvre that rule prohibits.
Three facts bear on whether it is admissible, and \S20.4 of the
pre-registration sets out the full argument for a reader who wishes to reject
it. The criterion is outcome-blind, depending only on a stratum's event count.
The direction is exculpatory: the reduction removed the family's only surviving
discovery, a language contrast resting on twelve events, rather than creating
one; a family narrowed to manufacture significance narrows \emph{around} its
significant results. And the verdict is invariant to the choice: no contrast in
E1 survives Benjamini--Hochberg under either denominator, the smallest adjusted
value being $q = \pcSubgroupMinQ$ at $m = \pcNSubgroupContrasts$.
Section~\ref{sec:equity_strata} reports the consequence.

The floor applies to E1 and \emph{not} to E2, and the asymmetry is deliberate
rather than an oversight. E1's estimand is an AUROC, which requires enough
events to order cases against non-cases and is not usefully estimable on a
handful; E2's is a stay-level proportion, which is estimable at any stratum size
and simply carries a wide interval when the stratum is small. E2 therefore
retains every stratum, including one insurance category resting on seven stays
whose interval spans nearly sixty percentage points and is reported with that
interval visible. Retaining it makes the correction \emph{stricter}, not more
permissive, because it enlarges the denominator over which false-discovery
control is applied; dropping it would be the manoeuvre that needs defending.

\subsection{Analysis Status Tagging}

Every result in Chapter~\ref{ch:evaluation} carries an explicit status tag
rather than relying on a single blanket statement in this chapter:

\begin{itemize}
    \item \textbf{Confirmatory}: estimand, threshold and direction were
          fixed in \texttt{00\_preregistration.md} before any modelling.
    \item \textbf{Exploratory}: pre-registered in scope but not in specific
          contrasts. Generates hypotheses; does not test them.
    \item \textbf{Descriptive}, no inferential claim is attached.
\end{itemize}

The tags are held in a machine-readable analysis register
(\texttt{12\_analysis\_\allowbreak register\allowbreak .parquet}) produced by
the pipeline, and are
reproduced in Table~\ref{tab:register}.

\subsection{Status of the Inference Plan Itself}

This inference and multiplicity plan is \textbf{Protocol Amendment 1}, dated
30 July 2026 and recorded as \S13--15 of
\texttt{00\_preregistration.md}. It was written \emph{after} the analyses of
Chapter~\ref{ch:evaluation} had been run and inspected, and it is therefore
itself post-hoc. It is reported as such rather than presented as
pre-specification.

The amendment changes no hypothesis, estimand, threshold, direction, subgroup
or model. It attaches uncertainty to contrasts that were already pre-specified
and imposes the multiplicity control the original protocol omitted. Where the
correction changes a verdict recorded in the original protocol, the corrected
verdict is the one reported, and the change is stated explicitly.

\subsection{Feature Ablation: Action-Derived Predictors (Protocol Amendment 5)}
\label{sec:ablation_method}

The central argument of this thesis is that the Sepsis-3 label is constituted
by clinician action rather than merely correlated with it. The same objection
applies, with equal force, to part of the \emph{feature} set, and it would be
inconsistent to press it against the outcome while leaving the predictors
unexamined.

Two kinds of predictor record clinician behaviour rather than patient
physiology. \textbf{Vasopressor exposure} is a treatment: it is something done
to the patient, not something measured in them. The \textbf{per-test order
indicators} (whether lactate, creatinine, bilirubin, platelets, white cells
or a blood gas was obtained in a given hour) record a decision to look. Their
predictive value is well established and is usually read as evidence that
measurement patterns encode clinical suspicion, which is precisely the problem:
a model that leans on them is partly learning when clinicians became worried.

\textbf{Protocol Amendment 5} (post-hoc, dated 6 August 2026 and recorded as
\S19 of \texttt{00\_preregistration.md}) therefore adds an ablated arm. The primary
model and the gradient-boosted comparator are re-fitted on the same training
rows, the same splits and the same hyperparameters, with the
\pcAblNDropped{} action-derived features removed and every physiological
measurement retained.

Three properties make the contrast interpretable:

\begin{itemize}
    \item \textbf{The ablation is deliberately conservative.} Forward-filled
          laboratory \emph{values} are retained even though a value exists only
          because someone ordered the test. Dropping them as well would remove
          most of the laboratory signal and confound the ablation with a loss
          of physiological information. What is removed is the set of features
          that encode \emph{only} clinician behaviour.
    \item \textbf{Nothing is re-tuned.} The gradient-boosted arm keeps the
          hyperparameters of the full fit. Re-tuning would confound the feature
          set with the search budget.
    \item \textbf{Both arms are scored by the same metric code on the same test
          rows}, so the difference is attributable to the feature set alone.
\end{itemize}

The arm is exploratory. It enters no hypothesis, no confirmatory family and no
variance decomposition; it is reported as estimation, and its two readings, the change in discrimination and the change in cross-label coefficient
stability, are given in Section~\ref{sec:ablation}.


\section{Pre-Registration and Reporting}

The full protocol was locked before any data extraction. All six hypotheses,
label variants, model classes, splits, metrics, and decision thresholds are
specified in \texttt{00\_preregistration.md}. Any analysis not specified there
is labelled as exploratory, and, following
Section~\ref{sec:inference}, each individual result carries its status tag in
the results table itself rather than inheriting it from this statement alone.
Reporting follows TRIPOD+AI \cite{collins2024tripodai}; code availability and
the constraints the data use agreement places on it are set out in
Section~\ref{sec:code_availability}.
